\chapter*{Tóm tắt}
\addcontentsline{toc}{chapter}{Tóm tắt}

\textbf{Tóm tắt:} Nhà thông minh dùng cảm biến môi trường có thể học thói quen sinh hoạt của người ở mà
không cần camera hay thiết bị đeo --- nền tảng cho các ứng dụng chăm sóc người cao tuổi như nhắc uống
thuốc. Tuy nhiên, khi trong nhà có hai người, cảm biến không biết ai đã kích hoạt nó, nên các phương
pháp tìm thói quen dành cho nhà một người trộn lẫn thói quen của hai người. Khóa luận nghiên cứu bài toán
\emph{học thói quen của từng cư dân} từ luồng sự kiện cảm biến trộn lẫn. Khóa luận phát biểu hình thức
khái niệm thói quen theo từng người, hiện thực nó bằng thuật toán quét cửa sổ đếm theo ngày, xây dựng bộ
đánh giá ở mức thói quen, và so sánh năm cách tiếp cận --- hai baseline không nhãn (trong đó có phương
pháp của đề tài thực tập), một baseline tra bảng, Random Forest và Bi-GRU --- trên năm nhà của bộ dữ liệu
CASAS, đánh giá theo thứ tự thời gian. Trên ba nhà hai người, nhận dạng (người, hoạt động) rồi áp định
nghĩa thói quen đạt F1 thói quen 0,50--0,71 (Random Forest), so với 0,22--0,33 của phương pháp theo vùng;
ở nhà Cairo, mô hình tìm đủ 6/6 thói quen gắn với từng người, kể cả thói quen uống thuốc vốn trùng giờ và
trùng cảm biến với bữa sáng. Phân tích lỗi lan truyền bằng cách làm nhiễu nhãn có kiểm soát cho thấy thói
quen chịu được lệch giờ nhỏ và nhận thừa ngẫu nhiên, nhưng nhạy với bỏ sót, nhầm người và lỗi có hệ
thống. Thí nghiệm phát hiện ngày bỏ thuốc trên dữ liệu có cài bất thường cho thấy Bi-GRU phát hiện được
9/10 ngày nhưng còn nhiều báo động nhầm, trong khi cách tiếp cận theo vùng không phát hiện được ngày nào.

\textbf{Từ khóa:} nhà thông minh, nhiều cư dân, cảm biến môi trường, khai phá thói quen, nhận dạng hoạt
động, phát hiện bất thường.

\vspace{1cm}
\begin{otherlanguage}{english}
\noindent\textbf{Abstract:} Smart homes with ambient sensors can learn the daily habits of their
residents without cameras or wearables --- a foundation for elderly-care applications such as medication
reminders. When two people share a home, however, the sensors do not know who triggered them, and habit
discovery methods designed for single-resident homes mix the habits of both. This thesis studies how to
\emph{learn per-resident habits} from such a mixed sensor stream. It formally defines a per-resident habit,
implements the definition as a day-counting window-scan algorithm, builds a habit-level evaluation
protocol, and compares five approaches --- two label-free baselines (including the author's internship
method), a lookup baseline, Random Forest and Bi-GRU --- on five CASAS homes under a chronological split.
On the three two-resident homes, recognising (resident, activity) and then applying the habit definition
reaches a habit F1 of 0.50--0.71 (Random Forest) versus 0.22--0.33 for the zone-based approach; in the
Cairo home it recovers all 6/6 per-resident habits, including a medication habit that shares its time and
sensors with breakfast. A controlled label-corruption analysis shows that habits tolerate small time
shifts and random false positives but are sensitive to missed occurrences, resident swaps and systematic
errors. In a missed-medication detection experiment with injected anomalies, Bi-GRU detects 9 of 10
skipped days at the cost of frequent false alarms, while the zone-based approach detects none.

\noindent\textbf{Keywords:} smart home, multi-resident, ambient sensors, habit mining, activity
recognition, anomaly detection.
\end{otherlanguage}
